% !PRINTED 188
\weberpage{188}

If, following \S.\,16, (4), we arrange the function
\[
F(x_1 + t\xi_1, x_2 + t\xi_2 \ldots x_m + t\xi_m)
\]
in powers of $t$, we obtain for the coefficient of $t^\nu$
\weq{4}{F_\nu(x,\xi) = \frac{1}{\Pi(\nu)}\sum^{\alpha}\frac{\Pi(\nu)}{\Pi(\alpha_1) \ldots \Pi(\alpha_m)}\xi_1^{\alpha_1}\xi_2^{\alpha_2}\ldots\xi_m^{\alpha_m}D_{\alpha_1,\alpha_2\ldots\alpha_m}(F),}
\[
\alpha_1 + \alpha_2 + \cdots + \alpha_m = \nu.
\]
If we apply the linear transformation (1), \S.\,55 to the variables
$\xi$ and $x$ simultaneously, then $F_\nu(x,\xi)$ passes over into $\Phi_\nu(x',\xi')$,
which is obtained from $F_\nu$ by simultaneously replacing all the letters
$x, \xi, a$ by $x', \xi', a'$ and $F$ by $\Phi$.

From this it follows:

I. \gesperrt{If one regards $F_\nu(x,\xi)$ as a form of degree $\nu$
in $\xi$ and forms an invariant of this form, whose coefficients therefore
still depend on the $x$, one obtains a covariant of the form $F$.}

And since the variables $x$ and $\xi$ are subject to the same
transformation, it likewise follows:

II. \gesperrt{If one forms a covariant of the function $F_\nu(x,\xi)$
regarded as a function of $\xi$, and in it replaces the $\xi$ by the $x$,
one obtains a covariant of the original function $F(x)$.}

The same theorems hold also for simultaneous invariants
and covariants.

For the functions $F_\nu(x,\xi)$ we have the theorem
\weq{5}{F_\nu(x,\xi) = F_{n-\nu}(\xi,x),}
which follows by comparing corresponding powers of $t$
on the two sides of the following identity:
\[
F(x_1+t\xi_1\ldots x_m+t\xi_m) = t^n F(t^{-1}x_1+\xi_1\ldots t^{-1}x_m+\xi_m)
\]

The functions
\[
P_\nu(x,\xi) = \frac{\Pi(\nu)\;\Pi(n-\nu)}{\Pi(n)}F_\nu(x,\xi)
\]
are called the \gesperrt{polars of the function $F(x)$}, and specifically
$P_\nu(x,\xi)$ is called the $\nu$th polar. For them too the following theorem holds:
\[
P_\nu(x,\xi) = P_{n-\nu}(\xi,x).
\]
